% !TEX program = xelatex
% BiT-128 恢复路线报告 — 编译: xelatex 两次
\documentclass[UTF8,11pt,a4paper]{ctexart}

\usepackage{amsmath,amssymb,bm}
\usepackage{booktabs,array,tabularx,longtable}
\usepackage[margin=2.3cm]{geometry}
\usepackage{enumitem}
\usepackage{xcolor}
\setlist{itemsep=2pt,topsep=4pt}
\usepackage[colorlinks=true,linkcolor=blue!55!black,urlcolor=blue!55!black,citecolor=blue!55!black]{hyperref}
\usepackage{fancyhdr}

\setmainfont{Cambria}
\setmonofont{Consolas}
\newcommand{\code}[1]{\texttt{\small #1}}

\pagestyle{fancy}\fancyhf{}
\fancyhead[L]{\small BiT-128 公开数据等价私钥恢复}
\fancyhead[R]{\small\thepage}
\renewcommand{\headrulewidth}{0.3pt}

\title{\vspace{-1.2cm}\textbf{BiT-128 公开数据等价私钥恢复：\\hint 通道的算子求逆}}
\author{GLM-5.3 Max \and JY Liu}
\date{2026 年 10 月 19 日}

\begin{document}
\maketitle
\thispagestyle{fancy}

\begin{abstract}\noindent
BiT-128 是 NICCS 首轮征集的签名候选（编号 sign-02），属 BLISS 一族的双模高斯设计。本文
报告一项仅使用公开数据的密钥恢复：输入为公钥与 400 万条普通签名，不借助侧信道、插桩或
计时信息，我们精确恢复了秘密多项式 \code{s0} 的全部 768 个系数，并按验证等价关系补全
$(e, b_0)$，构造出通过未修改参考 C 验证器的等价签名密钥——对新鲜消息的伪造签名 20/20
被接受，四组阴性对照全部被拒绝。恢复的核心是一个此前未被描述的泄露通道：签名侧 hint
论证式中的掩码消去结构，使得公开的量化 hint 成为等价私钥未知量 $R = b_0 - e$ 的直接
一阶观测；对该观测的确定性结构建立线性算子定律并求逆，即得 $R$。标签通道准确率
96.19\%，在 $4\times10^6$ 签名规模上直接分类零错误。同一密钥经两条相互独立的路线
（本文 hint 算子路线与 $z_1$ 回归路线）恢复出逐字节一致的等价私钥。

\medskip\noindent\textbf{Abstract.}
BiT-128 is a bimodal-Gaussian lattice signature candidate (submission sign-02 of
the NICCS first round). We report a public-data-only key recovery: from the
public key and 4 million ordinary signatures---no side channels, no
instrumentation, no timing---we recover all 768 coefficients of the secret
polynomial \code{s0} exactly, complete $(e,b_0)$ up to verification
equivalence, and forge signatures on fresh messages that pass the unmodified
reference C verifier 20/20, while all four negative controls are rejected. The
core of the attack is a previously undescribed leakage channel: in the signer's
hint identity, fresh masking noise cancels, so the public quantized hint is a
direct first-order observation of $R=b_0-e$, the only unknown of the equivalent
key. A linear operator law for the deterministic structure of this observation,
once inverted, yields $R$.
\end{abstract}

\section{对象、授权与结论边界}

审计于 2026 年 9 月 20 日启动，在 BiT 提交方授权下针对其自有参考实现进行；攻击输入仅为
公钥与公开签名，秘密密钥材料只在事后评分时使用。缺陷属设计层面——双模符号位的联合
分布携带密钥信息，而方案的安全论证只覆盖了边缘分布的独立性（详见技术报告的证明评估
一节，\code{paper/security\_proof\_assessment\_v2.pdf}）。

先给结论的量化形态，避免歧义：

\begin{table}[htbp]\centering\small
\begin{tabular}{@{}ll@{}}
\toprule
项 & 结果 \\
\midrule
输入 & 公钥 $+$ $4\times10^6$ 条公开签名（普通消息流） \\
\code{s0} & 768/768 系数精确，模 $q$ 解在窗口约束下唯一 \\
$(e,b_0)$ & 验证等价补全：公钥只确定 $R=b_0-e$，任何满足约束的组合签名能力相同 \\
伪造 & 新鲜消息，未修改参考验证器，20/20 接受 \\
阴性对照 & 硬币标签臂、篡改 $e$、篡改签名、篡改 pk，四组全部拒绝 \\
交叉验证 & 两条独立路线恢复同一把等价私钥（逐字节一致）；$3\times10^6$ 条签名经原始 C 验证器重验证零失败 \\
\bottomrule
\end{tabular}
\end{table}

\section{缺陷定位：联合分布对边缘分布}

BiT-128 的签名响应使用共享的双模符号位（BLISS 式 bimodal 结构）：秘密扰动的符号选择在
多个坐标间共享，验证方程按双模恒等式吸收符号。单个观察者从一条签名里看到的每个分量
都是边缘分布意义下``正确''的；问题出在\textbf{联合}分布——共享符号位使 $z_0$ 的取值
与秘密之间建立了公开可计算的弱相关。$z_0$ 通道由此给出 96.19\% 准确率的符号标签：不
完美，但远高于一半，足以驱动后续线性估计。

这一缺陷形态与 Espitau 等人对 BLISS 的侧信道攻击~[EFGT17] 有本质区别：后者需要运行时
泄露，本攻击完全被动、只读签名文本。

\section{恢复路线}

\subsection{hint 通道：掩码在哪里消去}

签名侧计算 hint 论证式
\[
w_a \;=\; w - z_2 + b_0\cdot c_{\mathrm{signed}}\quad(\text{含行 0 的 } b \text{ 修正}),
\]
其中的关键结构是：$w$ 与 $z_2$ 都含有新鲜掩码 $y_2$，在差 $w - w_a$ 中\textbf{相互
抵消}。于是
\[
w_a \;=\; \mathrm{base} - (R * c_{\mathrm{signed}}) + \cdots,\qquad R = b_0 - e,
\]
与新鲜掩码无关；公开的量化 hint $h$ 量化的是 $y_2 - (R * c_{\mathrm{signed}}) - \cdots$。
$y_2$ 是密度已知的新鲜噪声，$R$ 是等价私钥的全部未知量——$h$ 是 $R$ 的直接一阶观测。
这一消去结构是本攻击的核心观察，也是 BiT 的设计者未计入的泄露面。

\subsection{算子定律}

对带符号挑战的对齐统计量，其期望中的确定性结构是 $R$ 的线性泛函。对齐组合学部分在
``挑战位置不重复''的抽样模型下有精确期望，实际由 Monte-Carlo 校准为算子 $F_{\mathrm{op}}$；
接受条件化项没有解析模型，由实证界定——求逆后残差 $\sigma \approx 0.10$
（$4\times10^6$ 签名）。估计方程的形态：
\[
\bigl(M\cdot I + N\cdot F_{\mathrm{op}}\bigr)\,\hat R
\;=\; -\frac{\gamma_2}{s}\,\Phi \;-\; P_0\cdot M\cdot\Bigl(1-\frac{2m}{d}\Bigr)\cdot\mathbf{1}[i=0],
\]
标签尺度 $s$ 逐行以整数聚类目标拟合。

需要说明两处近似及其代价。其一，解析 tent 核是 iid 位置近似，对角超出精确 $N\tau$ 达
$(\tau-1)/d = 11.3\%$；相对位置互斥的 Monte-Carlo 算子把 Frobenius 误差压到 10.9\%。
其二，接受条件化未建模，这是本路线明确的诚实边界，其影响由实证残差界定而非声称排除。

\subsection{全解一致性修正}

解出 $\hat R$ 后，公钥关系 $s_0 = A_0^{-1}(16\,b_1 + R)$ 给出修正入口。我们采用一个
自设的验收准则（razor）：模 $q$ 解线性系统得到候选修正，\textbf{仅当整个 768 维解全部
落入 $\{-1,0,1\}$}（违规计数 $\mathrm{nbad}=0$）才接受；搜索按单点$\to$成对$\to$三元
阶段化推进。实际在 $N=4\times10^6$ 时直接分类已零错误，修正器是纵深防御而非必需环节。

\subsection{数据量阶梯}

\begin{table}[htbp]\centering\small
\begin{tabular}{@{}lll@{}}
\toprule
规模 & 直接分类准确率 & 说明 \\
\midrule
$10^6$（oracle 臂） & 99.87\% & 单修正闭合 \\
$10^6$（attack 臂） & 99.09\% & 7 个 iid 跨界 \\
$4\times10^6$ & 100\% & $s_0$ 768/768 精确，无需修正 \\
\bottomrule
\end{tabular}
\end{table}

\section{恢复程度的精确语义}

哪些量是精确恢复、哪些是等价补全，需要严格区分。$s_0$ 精确：与原钥完全一致，且经
两条独立路线交叉验证。$(e,b_0)$ 是等价补全：公钥只确定了 $R = b_0 - e$，任何满足
$\hat e \in \{-1,0,1\}$、$\hat b_0 = R + \hat e \in [-8,8)$ 的组合具有相同签名能力；
我们的补全在允许处一律取 $\hat e = 0$（26 个强制边界坐标）。验证器无法区分两种密钥的
签名——等价的含义是验证等价（不可区分签名能力），不声称输出同分布。

\section{与其它路线的关系}

针对 BiT-128 的密钥恢复目前共有三条独立路线：本仓库的 hint 算子路线、$z_1$ 回归
路线，以及 ePrint 2026/2232 第 8 章的 SSD 去卷积路线。三者恢复同一密钥，机制各异；
逐维度对比见 \code{paper/three\_way\_comparison.md}。本路线的独立价值在于 hint 通道的
显式代数刻画（掩码消去恒等式）与算子求逆框架——它是通用恢复框架 umodel（本仓库顶层
目录）的第一个已验证实例，在与 BiT 代数不同的循环环玩具实例上零改动完成了端到端
恢复，对 CS-128 与 COMPASS-SIG-128 的适用性判定也是同一框架给出的（见对应目录的报告）。

\section{下一步工作}

按代价从小到大：
\begin{enumerate}
\item \textbf{接受条件化项的解析建模}。当前是整条路线唯一未建模的环节，实证残差
  $\sigma \approx 0.10$。若能给出闭式或半经验修正，数据量需求有望从 $4\times10^6$
  显著下降（$z_0$ 标签在 $10^6$ 已 99\%+，瓶颈在算子精度）。
\item \textbf{低数据量体制}。$4\times10^6$ 是``直接零错误''的充分条件而非必要条件；
  修正器 $+$ 逐行标度拟合在 $10^6$ 已能闭合（99.87\%），系统性的数据量—精度曲线
  （$\sigma \propto N^{-1/2}$，实测拟合指数 $-0.497$）已在框架中建立，可用于外推更小
  $N$ 下的修正搜索代价。
\item \textbf{标签与 hint 通道的信息论最优融合}。两通道信息目前串行使用；联合似然的
  最优加权理论上还能压低阈值。
\item \textbf{向同族方案迁移}。凡签名侧存在``差式中新鲜掩码消去''结构的双模设计都可能
  暴露同类通道；umodel 的 SchemeSpec 声明式规范就是为这类迁移准备的（带符号挑战的矩
  算子已推导：$\mathbb{E}[C^{\mathsf T}C] = \tau I$ 精确成立，交叉项消失）。
\end{enumerate}

\section{复现指南}

\begin{verbatim}
# 路线 A：仅凭本目录公开工件（数分钟，纯 CPU）
python3 code/run_pipeline.py --artifacts artifacts --out key_hat.bin
# 预期输出 sha256:
# 4191a44b4bf1e081175986302b0d30605bb86f827dbe02c80d90ef7f704d5a11

# 路线 B：从零复现（需要 BiT 参考实现与约 400 万条签名）
python3 code/hint_channel_fast.py <dump.bin> <nsig> <tag>   # 累加 Phi_R
python3 code/hint_perrow.py real sim3m                      # 闭合
./forge_all pk.bin key_hat.bin 20    # 伪造闭环（原始 C 签名器/验证器）

# 通用框架自测与演示
python3 ../umodel/selftest.py
python3 ../umodel/run_demos.py
\end{verbatim}

秘密数据（真实密钥 \code{ground.bin}、\code{sk.bin}、签名 dump）不入库；\code{--score}
可选接受真值文件，仅用于事后准确率报告。八项交叉验证与全部带日期工件的清单见技术报告
可复现性一节。

\section{更正记录}

我们早期的笔记曾声称 $z_1$ 回归路线``精度饱和于 65--69\%、永不收敛''。该说法已撤回：
饱和是早期统计量 $\psi(z) = \mathrm{sign}(z)/(1024-|z|)$ 加符号阈值判决的伪象，对原始
响应的普通加权最小二乘按 $\mathrm{RMSE} \propto N^{-1/2}$ 收敛（RMSE 0.184 @
$2\times10^5$，经公钥修正后闭合）。如实记录于此，因为它说明了统计量构造对结论的
敏感性——这一教训同样适用于 CS-128 与 COMPASS-SIG-128 的判定（见各自报告）。

\section*{参考文献}
\begin{enumerate}[label={[\arabic*]}]
\item L.~Ducas, A.~Durmus, T.~Lepoint, V.~Lyubashevsky. \emph{Lattice Signatures and
  Bimodal Gaussians}. CRYPTO 2013（ePrint 2013/383）.
\item V.~Lyubashevsky. \emph{Fiat--Shamir with Aborts}. ASIACRYPT 2009.
\item T.~Espitau, P.-A.~Fouque, B.~G\'erard, M.~Tibouchi. \emph{Side-Channel Attacks on
  BLISS}. CCS 2017（ePrint 2017/505）——需要运行时泄露；本攻击不需要.
\item Z.~Xiong, M.~Wang. \emph{Design-Level Cryptanalysis of the ICCS NGCC Round-1
  Public-Key Candidates}. ePrint 2026/2232（SSD 去卷积路线）.
\item BiT 提交包，NICCS 2026 首轮（sign-02）.
\end{enumerate}

\end{document}
